\input{_actp.latex}

% --- Titre et classe
\def\Titre{... La cagnotte ?}
\def\Classe{2dne}

% ---
\begin{document}
\vspace*{0.1cm}
\MonTitre{\ldots}
\vspace*{0.4cm}
\section{Lettre n\up{o}1}

\img{7cm}{squaro_R.png}

Dans chacune des cases se trouve un nombre, de $0$ à $4$ : il correspond au nombre de ronds à colorier parmi ceux situés aux quatre coins de cette case.

\vspace*{0.4cm}

\section{Lettre n\up{o}2}
Une série de $5$ valeurs a une moyenne de $10$.
Quatre valeurs sont connues : $7$ ; $9$ ; $13$ ; $17$.

La cinquième valeur manquante est une lettre de l'alphabet ($A=1$, $B=2$, \ldots).  Quelle est-elle ?

\vspace*{0.4cm}

\section{Lettre n\up{o}3}
Dans un donjon, un coffre est scellé par un sort.

L'inscription gravée sur la pierre dit~:

\medskip
\og~Lance deux dés à 6 faces. Compte toutes les façons d'obtenir
une somme \textbf{strictement supérieure à 8}, puis multiplie
ce nombre par $2$ et ajoute $2$.

La lettre du code est à cette position
dans l'alphabet~($A=1$, $B=2$, \ldots).~\fg

\medskip
Quelle lettre ouvre le coffre~?

\vspace*{0.4cm}

\section{Lettre n\up{o}4}

Un hacker a chiffré un message~: chaque lettre du code correspond
à un point dans un repère. Une seule lettre est \og dans la combine \fg~:
c'est celle dont les coordonnées vérifient l'équation de la droite
$\mathcal{D}$~:
\[
	y = 2x - 1
\]

\medskip
\begin{center}
	\begin{tabular}{|c|c|c|c|c|c|c|c|}
		\hline
		\textbf{A}  & \textbf{B}  & \textbf{C}  & \textbf{D}
		            & \textbf{E}  & \textbf{F}  & \textbf{G}   & \textbf{H}   \\
		\hline
		$(0\,;\,0)$ & $(1\,;\,3)$ & $(2\,;\,2)$ & $(-1\,;\,1)$
		            & $(3\,;\,4)$ & $(0\,;\,2)$ & $(3\,;\,5)$  & $(-2\,;\,1)$ \\
		\hline
	\end{tabular}
\end{center}

\medskip
Quel point est sur $\mathcal{D}$~? C'est la lettre du code.

\vspace*{0.5cm}

\section{ ... et le dernier (nombre)}

Pour vaincre le boss final, tu dois lancer \textbf{4 sorts}
(Feu, Glace, Foudre, Ombre) dans un certain ordre.
Toutes les séquences sont possibles\ldots\ sauf deux~:

la séquence Feu~$\to$ Glace~$\to$ Foudre~$\to$ Ombre et son exact miroir, car elles déclenchent une explosion fatale.

Combien de séquences utilisables reste-t-il~?

\vspace*{2cm}

\centering\Huge{\ldots \quad \ldots \quad \ldots \quad en \quad  \ldots \quad \ldots}

\end{document}
